\begin{lemma}[Controlled incident colors persist through a series interval]
Let \(\mathcal T\) be a certified Tashkinov series with \(r\)
noninitial rounds.  Let \(p,q<r\), with \(p\leq q\), and put
\[
 i=p+1,\qquad j=q+1,\qquad T_i=\mathcal T.\mathrm{tree}(p+1).
\]
Suppose that the exact edge identity \(f\) is incident with \(T_i\),
that \(\varphi_{i-1}(f)=c\), and that
\[
 c\in\overline\varphi_{i-1}(T_i)\cup D_{i-1}.
\]
Then either \(\varphi_j(f)=\varphi_{i-1}(f)\), or there is a round
\(h\) with \(i\leq h\leq j\) for which
\(\Theta_h=\mathrm{PE}\) and the exact singleton connecting-edge set is
\(F_h=\{f\}\).  In particular, if
\(f\in E(G[T_i])\), then
\[
 \varphi_j(f)=\varphi_{i-1}(f).
\]
Here paper round \(h\) is represented by the positive data slot
\((h-1).\mathrm{succ}\).
\end{lemma}
\begin{proof}
We induct over the finite ordered interval of zero-based rounds
\(p,p+1,\ldots,q\).  The induction invariant after round \(a+1\) has
three parts.  First, \(f\) is incident with \(T_{a+1}\).  Second, its
current value has the form \(\varphi_{a+1}(f)=d_a\), where
\[
 d_a\in\overline\varphi_{a+1}(T_{a+1})\cup D_{a+1}.              \tag{1}
\]
Third, either its current value is its initial value
\(\varphi_p(f)\), or one of the already processed rounds is a PE whose
singleton connecting-edge set is \(\{f\}\).

The incidence part persists because every noninitial constructor in
\noderef{TashkinovSeriesStep} declares the preceding tree to be an
ordered initial segment of the next one.  Expanding
\noderef{TreeSequenceInitialSegment}, every old tree vertex remains a
new tree vertex, so the witness in \noderef{TreeIncidentEdge} remains a
valid incidence witness.

For the first round, apply \noderef{HkuIncidentColorPreserved} to
\(f,c\) and the hypotheses of the lemma.  Its first conclusion says
that the value is preserved or that this exact round is a PE with
singleton connecting-edge set \(\{f\}\).  Its second conclusion supplies
the color \(d_p\) and membership (1).  This establishes the invariant.

Now suppose the invariant holds through round \(a+1<q+1\).  The second
inclusion in \noderef{HkuMissingStoredMonotone} enlarges (1), under the
fixed coloring \(\varphi_{a+1}\), from \(T_{a+1}\) to \(T_{a+2}\).
Thus the current exact color satisfies the controlling-set hypothesis for
the next application of \noderef{HkuIncidentColorPreserved}.  That
application gives both the next fixed-tree membership and the following
dichotomy: the exact edge value is unchanged, or the new round is a PE
whose singleton connecting-edge set is \(\{f\}\).  The first inclusion
of \noderef{HkuMissingStoredMonotone}, which is built into that one-step
conclusion, accounts for the change from the old coloring to the new one;
its second inclusion will again prepare the following round.  In the
unchanged case the old disjunction in the invariant is retained.  In the
exceptional case the new round itself witnesses that disjunction.  This
completes the interval induction.  At its last index it gives precisely
the first assertion of the lemma.

Assume now that \(f\in E(G[T_i])\).  By the definition
\noderef{InternalEdges}, every endpoint incident with \(f\) belongs to
\(V(T_i)\).  The vertex containment obtained above from
\noderef{TreeSequenceInitialSegment} therefore makes \(f\) internal to
every \(T_h\) with \(i\leq h\leq j\).

It remains to exclude the alternative produced by the induction.  In a
PE constructor, the singleton stored in \(F_h\) is the exact edge
identity selected by the \noderef{MaximumDefectiveChoice}.  Its
\(\mathrm{edgeDefective}\) field and the definition
\noderef{DefectiveEdge} put that identity in the corresponding
\noderef{BoundaryColorEdges}, hence in the exact boundary set
\noderef{BoundaryEdges} of \(T_h\).  Equality of singleton finite sets
would therefore make \(f\) a boundary edge of \(T_h\).  A boundary edge
has an endpoint outside \(V(T_h)\), while internality says that every
endpoint of the same exact identity lies inside \(V(T_h)\), a
contradiction.  Thus no PE exception is possible, and the final color is
the initial color as claimed.
\end{proof}
